\chapter{Bewertung und Lessons Learned}

Das Projekt beweist, dass mit überschaubaren Mitteln und Freiverfügbarkeit von Komponenten wie Raspberry Pi, Spionspiegel und Kamera funktionale Smart Mirror-Prototypen umgesetzt werden können. Die enge Verknüpfung von Hardware und Software erforderte eine iterative Herangehensweise, wie sie im modernen Projektmanagement (agil/iterativ) beschrieben wird. Der Defekt am Raspberry zeigte, wie wichtig zeitliche Puffer und Ersatzteile für eine fortlaufende Entwicklung sind.

Erfolgreich war insbesondere der modulare Aufbau, der schnelle Isolierung und Lösung von Teilproblemen erlaubte. Die zentrale Anforderung, die aufbereitete Anzeige von Vitalparametern, wurde mit guten Ergebnissen erfüllt.